两个同心理想导体球面半径分别为 $a,b$（$0<a<b$），内球带自由电荷 $+Q$，外球带 $-Q$，$Q>0$。区域 $a<r<c$ 填充介电常数为 $\epsilon_1>\epsilon_0$ 的各向同性均匀线性电介质，区域 $c<r<b$ 填充介电常数为 $\epsilon_2>\epsilon_0$ 的另一介质，$a<c<b$；界面没有自由面电荷，介质均无自由体电荷。
(1) 由电位移高斯定律求两个介质区域的径向场及球面间电势差；(2) 求总电容与静电能，并通过场能积分独立检验；(3) 求 $r=a$、$r=c$ 和 $r=b$ 上所有介质束缚面电荷密度，验证合计的束缚电荷为零；(4) 检验 $\epsilon_1=\epsilon_2$ 的均匀介质极限以及 $c\to a^+$ 的极限，解释为何界面法向电场一般不连续。
